\section{Circulant Matrices}

{\bf 1. The Structure of a Circulant Matrix}
An \(n \times n\) circulant matrix \(C\) is completely defined by its first row vectors  \((c_0, c_1, \dots, c_{n-1})\). 
Each subsequent row is a right circular shift of the preceding row:
\[C=\left(\begin{matrix}c_{0}&c_{1}&c_{2}&\dots &c_{n-1}\\ c_{n-1}&c_{0}&c_{1}&\dots &c_{n-2}\\ c_{n-2}&c_{n-1}&c_{0}&\dots &c_{n-3}\\ \vdots &\vdots &\vdots &\ddots &\vdots \\ c_{1}&c_{2}&c_{3}&\dots &c_{0}\end{matrix}\right)\]
We can express \(C\) as a polynomial combination of a fundamental cyclic shift matrix 
\(P\): \[P=\left(\begin{matrix}0&1&0&\dots &0\\ 0&0&1&\dots &0\\ \vdots &\vdots &\vdots &\ddots &\vdots \\ 0&0&0&\dots &1\\ 1&0&0&\dots &0\end{matrix}\right)\]
Because shifting \(n\) times returns the matrix back to the identity (\(P^n = I\)), the characteristic equation of this shift operator is \(P^n - I = 0\). 
Consequently, its eigenvalues must satisfy \(\lambda^n = 1\), which explicitly forces the \(n\)-th roots of unity into the equation.
Any circulant matrix can be built directly out of \(P\):\(C=c_{0}I+c_{1}P+c_{2}P^{2}+\dots +c_{n-1}P^{n-1}\) \\

{\bf 2. Roots of Unity as Eigenvectors}
Let \(\omega \) be the primitive \(n\)-th root of unity: \(\omega =e^{\frac{2\pi i}{n}} \).
The distinct \(n\)-th roots of unity are given by \(\omega^0, \omega^1, \omega^2, \dots, \omega^{n-1}\).
If we construct a test vector \(v_{k}\) using the powers of the \(k\)-th root of unity:
\[v_{k}=\left(\begin{matrix}1\\ \omega^{k}\\ \omega^{2k}\\ \vdots \\ \omega^{(n-1)k}\end{matrix}\right)\]
Multiplying the shift matrix \(P\) by this vector shifts every element up by one position, moving the bottom element to the top:
\[Pv_{k}=\left(\begin{matrix}\omega^{k}\\ \omega^{2k}\\ \vdots \\ \omega^{(n-1)k}\\ 1\end{matrix}\right)=\omega^{k}\left(\begin{matrix}1\\ \omega^{k}\\ \vdots \\ \omega^{(n-2)k}\\ \omega^{-k}\end{matrix}\right)\]
Because \(\omega^n = 1,\)  we know that \(\omega^{-k} = \omega^{(n-1)k}\). Substituting this back in reveals that \(v_{k}\) is an exact eigenvector of \(P\):\(Pv_{k}=\omega^{k}v_{k}\) \\

{\bf 3. Calculating the Eigenvalues of \(C\)}
Because \(C\) is a polynomial in terms of \(P\), it shares the exact same eigenvectors (\(v_{k}\)) as \(P\). 
We can calculate the corresponding eigenvalues (\(\lambda _{k}\)) of the circulant matrix by applying the polynomial equation directly to the roots:

\[ Cv_{k} = \left(\sum _{j=0}^{n-1}c_{j}P^{j}\right)v_{k} = \left(\sum _{j=0}^{n-1}c_{j}(\omega^{k})^{j}\right)v_{k}\] \[\lambda _{k} = c_{0}+c_{1}\omega^{k}+c_{2}\omega^{2k}+\dots +c_{n-1}\omega^{(n-1)k} \]

This structure shows that every eigenvalue \(\lambda _{k}\) of a circulant matrix is simply the evaluation of the row's coefficients mapped against the \(n\)-th roots of unity.\\

{\bf 4. Definition of the DFT Matrix}

An $n$-point DFT matrix, denoted by $F$, is built using the primitive $n$-th root of unity: $\omega  =e^{\frac{2\pi i}{n}}$.
The entry in the $j$-th row and $k$-th column (using 0-indexed notation from $0$ to $n-1$) is defined as: \[F_{j,k} = \frac{1}{\sqrt{n}}\, \omega^{j\cdot k} = \frac{1}{\sqrt{n}}\, e^{\frac{2\pi i}{n} j\cdot k }\]

By stacking all \(n\) eigenvectors together as columns, we form a unitary change-of-basis matrix \(F\), universally known as the Vandermonde / Discrete Fourier Transform (DFT) matrix:
\[F = \frac{1}{\sqrt{n}} 
\left(\begin{matrix}1&1&1&\dots &1\\ 
1&\omega &\omega^{2}&\dots &\omega^{n-1}\\ 
1&\omega^{2}&\omega^{4}&\dots &\omega^{2(n-1)}\\ 
\vdots &\vdots &\vdots &\ddots &\vdots \\ 
1&\omega^{n-1}&\omega^{2(n-1)}&\dots &\omega^{(n-1)(n-1)}
\end{matrix}\right)\] \\

{\bf 5. Properties of the DFT Matrix}

{\bf Symmetric}: 
The matrix is completely symmetric (\(F = F^T\)) because the index product $j\cdot k$ is commutative.\\

{\bf Orthogonal}: The rows and columns are completely orthogonal to each other.\\ 

{\bf Unitary}: We have  \(F^\dagger F = I\), where \(F^{\dag }\) is the conjugate transpose of \(F\). This means that
the conjugate transpose of $F$ is its inverse.\\ 

{\bf Invertible}: The inverse transform matrix (IDFT) simply changes the sign of the exponent (negating the imaginary components) and divides by $n$.\\

{\bf 6. Diagonalization via the DFT Matrix}
Using the $F$ matrix, any circulant matrix \(C\) can be completely diagonalized:
\[C=F^{\dag }\Lambda F\]
where  \[\Lambda = \text{diag}(\lambda_0, \lambda_1, \dots, \lambda_{n-1})\] is the diagonal matrix containing the eigenvalues computed from the roots of unity.\\

{\bf Eigenvalues and Determinant of a Circulant}
Any circulant matrix $C$ can be completely diagonalized using the Fourier matrix. If the first row of $C$ is given by the vector 
\[c = [c_0, c_1, \dots, c_{n-1}],\] 
its eigenvalues \(\lambda _{k}\) are simply the Discrete Fourier Transform (DFT) of the vector c:
\[\lambda _{k}=\sum _{j=0}^{n-1}c_{j} e^{\frac{2\pi i}{n} j\cdot k } \]
Because of this, the determinant of a circulant matrix is the product of all its eigenvalues:
\[\det (C)=\prod _{k=0}^{n-1}\lambda _{k}.\]

{\bf How Circulants Connect with Abel's method of Rationalization} Abel's method of rationalization is used to eliminate radical expressions from denominators of a ratio of polynomials $P(x)/Q(x)$ and relies on multiplying an algebraic 
$Q(x)$ expression by its ``conjugates'' $Q(\omega x), Q(\omega^2 x), \hdots, Q(\omega^{n-1} x)$ over a radical extension field, for $\omega$ a primitve $n$-th root of unity\\

{\it Roots of Unity}: In an \(n\times n\) circulant determinant, the eigenvalues of the matrix are given by evaluating a polynomial at the \(n\)-th primitive roots of unity (\(\omega^{k}\)).\\

{\it The Product Rule}: The determinant of a circulant matrix equals the product of all its eigenvalues:
\[\det (C)=\prod _{k=0}^{n-1}(c_{0}+c_{1}\omega^{k}+c_{2}\omega^{2k}+\dots +c_{n-1}\omega^{(n-1)k})\] 

{\it Rationalization}: In Abel's method, if you want to rationalize a denominator containing a radical like \(\sqrt[n]{a}\), you treat \(\sqrt[n]{a}\) as the variable \(x\) in a polynomial.
By mapping the denominator's coefficients into a circulant matrix, the resulting determinant calculates the product of the expression across all its field conjugates.
Because this product covers every root of unity, all the complex and irrational radical terms cancel out,
leaving a completely rationalized, pure integer or polynomial value as the final determinant. \\ 

{\bf Twisted Circulant Matrix}

A twisted Circulant matrix (or a Twistulant matrix) is a square $n \times n$ matrix where each row is shifted cyclically to the right by one position, much like a traditional circulant matrix. 
However, when the final element wraps around to the beginning, it is multiplied by a specific complex scalar $d$. \\

%{\bf Structure}
%For an $n \times n$ matrix defined by the first row \((c_0, c_1, c_2, \dots, c_{n-1})\) and a twist factor $\lambda$, 
%the matrix takes the following form:
%\[\left[\begin{matrix}c_{0}&c_{1}&c_{2}&\dots &c_{n-1}\\ 
%\lambda c_{n-1}&c_{0}&c_{1}&\dots &c_{n-2}\\ 
%\lambda c_{n-2}&\lambda c_{n-1}&c_{0}&\dots &c_{n-3}\\ 
%\vdots &\vdots &\vdots &\ddots &\vdots \\ 
%\lambda c_{1}&\lambda c_{2}&\lambda c_{3}&\dots &c_{0}\end{matrix}\right]\] \\

{\bf Summary of Rationalizing Denominators in Abel's Ratio $P(x)/Q(x)$}\\

{\bf 1. The General \(n \times n\) Matrix Definition}
Let \(\alpha = \sqrt[n]{d}\) be an algebraic radical of degree \(n\). 
The algebraic denominator expression to be cleared is defined as:
\(Q(\alpha )=\sum _{j=0}^{n-1}a_{j}\alpha ^{j}=a_{0}+a_{1}\alpha +a_{2}\alpha ^{2}+\dots +a_{n-1}\alpha ^{n-1}\)
The associated \(n \times n\) twisted circulant matrix \(C\) is structured as follows:
\[C=\left(\begin{matrix}a_{0}&a_{1}&a_{2}&\dots &a_{n-2}&a_{n-1}\\ d\cdot a_{n-1}&a_{0}&a_{1}&\dots &a_{n-3}&a_{n-2}\\ d\cdot a_{n-2}&d\cdot a_{n-1}&a_{0}&\dots &a_{n-4}&a_{n-3}\\ \vdots &\vdots &\vdots &\ddots &\vdots &\vdots \\ d\cdot a_{2}&d\cdot a_{3}&d\cdot a_{4}&\dots &a_{0}&a_{1}\\ d\cdot a_{1}&d\cdot a_{2}&d\cdot a_{3}&\dots &d\cdot a_{n-1}&a_{0}\end{matrix}\right)\] 

{\bf Eigenvectors}:
Let \(\omega = e^{\frac{2\pi i}{n}}\) be the primitive \(n\)-th root of unity.
The eigenvectors are independent of the specific matrix elements \(a_0, a_1, \dots, a_{n-1}\) and rely solely on \(n\) and \( \alpha \).
They can be explicitly defined by a scaled Discrete Fourier Transform (DFT) basis.
For \(k = 0, 1, \dots, n-1\), the \(k\)-th right eigenvector \(v_{k}\) of \(C\) is given by:
\[v_k = \begin{bmatrix}  1 \\  \omega^{-k} \alpha  \\  \omega^{-2k} \alpha^{2} \\  \vdots \\  \omega^{-(n-1)k} \alpha^{(n-1)}  \end{bmatrix}. \]

{\bf Eigenvalues:} While a standard circulant matrix uses pure Fourier modes as eigenvectors, a $d$-circulant matrix's eigenvectors account for the scalar shift $d$. 
Each eigenvalue \(\lambda _{k}\) corresponds to the following exact formula:  
\[ \lambda_k = \sum_{j=0}^{n-1} a_j \left(\omega^k \alpha \right)^j \] \\

{\bf Uncoupled Diagonalization}: Because all twisted circulant matrices commute with the twisted shift operator, they are simultaneously diagonalizable. \\

{\bf Standard Circulant}: When \(d = 1\), the matrix collapses into a standard circulant matrix, and the eigenvectors become the standard columns of the inverse DFT matrix, 
defined entirely by the roots of unity.\\

{\bf 2. The Determinant (Field Norm Formula)} \\

The determinant factors into the product of the field conjugates:
\[\det (C)= \prod_{k=0}^{n-1}\, \lambda_k  =  \prod _{k=0}^{n-1}\left(\sum _{j=0}^{n-1}a_{j} \left(\omega^{k} \alpha \right)^{j}\right)\]
Since the left-hand side involves the determinant of an $n \times n$ matrix whose elements are free of \(\omega \) and \(\alpha  =  \sqrt[n]{d} \), therefore expanding the product on the right-hand side eliminates all instances of \(\omega \) and \(\alpha  \).
This determinant is guaranteed to be a purely rational number if all \(a_{j}\) and $d$ are rational. 
It yields a purely rational polynomial given by the determinant of $C$:
\[\det (C)=\text{Norm}_{\mathbb{Q}(\sqrt[n]{d})/\mathbb{Q}}(Q(\alpha ))\in \mathbb{Q}\] 

{\bf 3. Clearing the Denominator via Classical Adjoint}
To rationalize the fraction \(\frac{1}{Q(\alpha )}\), use the matrix inversion property involving the Classical Adjoint matrix 
\(\text{adj}(C)\): \[C\cdot \text{adj}(C)=\det (C)\cdot I_{n}\]
The vector of algebraic basis elements is defined as:
\[\vec{\alpha }=\left(\begin{matrix}1\\ \alpha \\ \alpha ^{2}\\ \vdots \\ \alpha ^{n-1}\end{matrix}\right)\]
By the construction of the twisted circulant matrix, we have the identity:
\[C\vec{\alpha }=Q(\alpha )\vec{\alpha }\]
Multiplying both sides by the Classical Adjoint matrix 
\(\text{adj}(C)\) yields:\[\det (C)\vec{\alpha }=Q(\alpha )\cdot \text{adj}(C)\vec{\alpha }\]
Extracting the first component (the top row) gives the exact rationalization identity:
\[\frac{1}{a_{0}+a_{1}\alpha +\dots +a_{n-1}\alpha ^{n-1}}=\frac{\sum _{j=0}^{n-1}M_{0,j}\alpha ^{j}}{\det (C)}\]
Where \(M_{0,j}\) represents the signed \((0,j)\) cofactor of the first row of matrix \(C\):
\[M_{0,j}=(-1)^{j}\det (C_{\text{minor}(0,j)})\]